
\documentclass[12pt,letterpaper]{article}

% Paquetes esenciales
\usepackage[utf8]{inputenc}
\usepackage[spanish]{babel}
\usepackage{amsmath, amssymb, amsfonts}
\usepackage{graphicx}
\usepackage{float}
\usepackage{booktabs}
\usepackage[margin=1in]{geometry}
\usepackage{hyperref}

% Datos del laboratorio
\newcommand{\labnumero}{1}
\newcommand{\labtitulo}{Ley de Ohm y Kirchhoff}
\newcommand{\autor}{[Nombre del Estudiante]}
\newcommand{\fecha}{\today}

\begin{document}

% ========== PORTADA ==========
\begin{titlepage}
    \centering
    {\Large \textbf{Universidad Nacional Tecnológica de Lima Sur}}
    \vspace{0.5cm}
    
    {\large \textbf{Facultad de Ingeniería Electrónica}}
    \vspace{1cm}
    
    {\Huge \textbf{Informe de Laboratorio N°\labnumero}}
    \vspace{0.3cm}
    {\huge \textbf{\labtitulo}}
    \vspace{1.5cm}
    
    {\large \textbf{Estudiante:} \autor}
    \vspace{0.3cm}
    
    {\large \textbf{Fecha:} \fecha}
    \vspace{1cm}
    
    \rule{0.6\textwidth}{0.5pt}
    
    \vfill
    {\small \today}
\end{titlepage}

% ========== ÍNDICE ==========
\tableofcontents
\newpage

% ========== SECCIONES ==========
\section{Objetivos}
\begin{itemize}
    \item Objetivo general 1
    \item Objetivo específico 1
    \item Objetivo específico 2
\end{itemize}

\section{Marco Teórico}
\subsection{Ley de Ohm}
\begin{equation}
    V = I \cdot R
    \label{eq:ohm}
\end{equation}

\subsection{Leyes de Kirchhoff}
\begin{equation}
    \sum V = 0 \quad \text{(LVK)}
\end{equation}
\begin{equation}
    \sum I = 0 \quad \text{(LCK)}
\end{equation}

\section{Materiales y Equipos}
\begin{itemize}
    \item Fuente de voltaje DC (0-30V)
    \item Multímetro digital
    \item Resistencias: 100Ω, 220Ω, 330Ω, 470Ω, 1kΩ
    \item Protoboard
    \item Cables de conexión
\end{itemize}

\section{Procedimiento}
\begin{enumerate}
    \item Medir el valor real de cada resistencia
    \item Armar el circuito de la Figura \ref{fig:circuito}
    \item Aplicar voltajes de 1V a 10V
    \item Registrar mediciones
\end{enumerate}

\section{Resultados}
\begin{table}[H]
    \centering
    \caption{Mediciones realizadas}
    \begin{tabular}{|c|c|c|}
        \hline
        \textbf{V entrada (V)} & \textbf{I medida (mA)} & \textbf{R calculada (Ω)} \\
        \hline
        1.0 & 9.8 & 102.0 \\
        3.0 & 29.5 & 101.7 \\
        5.0 & 49.2 & 101.6 \\
        7.0 & 68.9 & 101.6 \\
        9.0 & 88.6 & 101.6 \\
        \hline
    \end{tabular}
    \label{tab:mediciones}
\end{table}

\section{Conclusiones}
\begin{itemize}
    \item Conclusión 1
    \item Conclusión 2
\end{itemize}

\end{document}
